\section{The combinatorial reduction arc (UC1--UC8)}
\label{sec:combinatorial}

This section consolidates the program's native combinatorial attack. It
is a chain of \emph{reductions}: each note replaces the residual
difficulty of its predecessor by a strictly smaller, more explicit
object, until Frankl becomes a single finitary integer inequality. The
arc proves real theorems --- it settles Frankl on a growing hierarchy
of structured sub-families and recovers the classical small-generator
cases natively --- but it does \emph{not} prove the conjecture. We say
so precisely in Section~\ref{ssec:combinatorial-status}.

Throughout this section $\F$ is union-closed with $U=\bigcup\F\in\F$.

\subsection{The compatibility-geometry framework (UC1)}
\label{ssec:uc1-uc2}

The founding manifesto (UC1) is a feasibility study. It checks whether
the two proof engines used on the $1/3$--$2/3$ conjecture --- a
Bousfield--Kan / Cheeger-expansion engine and a cohomological
moduli-space engine --- transfer to Frankl. Its honest verdict is
\emph{partial}: the structural framework transfers cleanly (the
dictionary of Section~\ref{ssec:analogy}; the transport graph $G(\F)$;
the Walsh picture), but \emph{neither engine} transfers. The transport
graph $G(\F)$ is disconnected in general (already
$\F=\{\emptyset,\{1,2\}\}$), so there is no irreducible transport chain
to run an expansion argument on; and the natural moduli space of all
union-closed families on $[n]$ is doubly exponential, with no
homology-sphere structure, so the moduli engine has nothing to act on.

The manifesto's lasting contributions are three. First, it pins the
intrinsic object: $\F\subseteq\cube$ with its induced transport graph
$G(\F)$, with the level-$1$ Walsh spectrum carrying the obstruction's
content (Proposition~\ref{prop:four-forms}). Second, it isolates the
\emph{orientation} of union closure as the engine to build a native
argument around (Remark~\ref{rem:asymmetry}): $\F_x$ is a filter,
$\F_{\bar x}$ an ideal. Third, it records a genuine negative result
that disciplines everything downstream.

\begin{proposition}[Globalisation defect is orthogonal to the bias, UC2 \S4]
\label{prop:globalisation-orthogonal}
Let $\mathrm{Def}(\F)=n|\F|-2|E(G(\F))|$ be the total number of
coordinate-toggle pairs at which $\sigma_x$ fails to be defined on
$\F$. Then $\mathrm{Def}(\F)$ cancels identically out of $\sum_x
\bias_x$: one has
$\sum_x|\F^{\stuck}_x|=\sum_{A\in\F}|A|-|E(G(\F))|$ and
$\sum_x|\F^{\stuck}_{\bar x}|=n|\F|-\sum_{A\in\F}|A|-|E(G(\F))|$, and
the edge term cancels in the difference.
\end{proposition}

\begin{proof}
Count edges of $G(\F)$ by direction: $|M_x|$ direction-$x$ edges, and
$|M_x|=|\F_x|-|\F^{\stuck}_x|$. Sum over $x$ and use
$\sum_x|\F_x|=\sum_A|A|$, $\sum_x|\F_{\bar x}|=n|\F|-\sum_A|A|$. The
difference is $\sum_x\bias_x=2\sum_A|A|-n|\F|$, with $|E(G(\F))|$
absent.
\end{proof}

Proposition~\ref{prop:globalisation-orthogonal} rules out the naive
mechanism ``failure to globalise forces the bias'': the \emph{unsigned}
globalisation defect is provably independent of $\sum_x\bias_x$. A
native argument must therefore use the \emph{signed} / \emph{oriented}
defect --- the filter/ideal structure --- not the raw count. This is
the discipline the rest of the arc obeys.

\subsection{Transport-deficiency and the trace-partition theorem (UC2)}
\label{ssec:uc2}

Fix a $\subseteq$-minimal nonempty member $W\in\F$ and put
$k=|W|=m(\F):=\min\{|A|:A\in\F,A\neq\emptyset\}$. We call $W$ a
\emph{minimal generator}. Every $A\in\F$ splits as
$A=S\sqcup R$ with $S=A\cap W\subseteq W$ and $R=A\setminus W$.

\begin{definition}
\label{def:trace-class}
The \emph{trace-classes} of $\F$ relative to $W$ are
$\F^S=\{A\in\F:A\cap W=S\}$ for $S\subseteq W$; they partition $\F$
into $2^k$ blocks. Write the \emph{level-$1$ potential}
$\phi(S)=2|S|-k$ and, for a family $\G$ of subsets of $W$,
$\Phi(\G)=\sum_{S\in\G}\phi(S)$.
\end{definition}

\begin{theorem}[Trace-partition theorem, UC2 \S3]
\label{thm:trace-partition}
For any minimal generator $W$,
\[
  \sum_{x\in W}\bias_x \;=\; \sum_{S\subseteq W}|\F^S|\,\phi(S)
  \;=\;\sum_{S\subseteq W}|\F^S|\,(2|S|-k).
\]
Consequently:
\begin{enumerate}[label=\textup{(\arabic*)},leftmargin=2.4em]
\item the map $A'\mapsto A'\cup W$ is an injection
  $\F^\emptyset\hookrightarrow\F^W$;
\item if $k\le2$ then $\sum_{x\in W}\bias_x\ge0$, so some $x\in W$ is
  abundant;
\item if $k\ge3$ then
  $\sum_{x\in W}\bias_x=k\bigl(|\F^W|-|\F^\emptyset|\bigr)
   +\sum_{0<|S|<k}|\F^S|\,(2|S|-k)$.
\end{enumerate}
\end{theorem}

\begin{proof}
For $x\in W$, $\bias_x=\sum_S|\F^S|\,\varepsilon_x(S)$ with
$\varepsilon_x(S)=+1$ if $x\in S$ and $-1$ otherwise; summing over
$x\in W$ replaces $\varepsilon_x(S)$ by $|S|-(k-|S|)=\phi(S)$. (1):
$A'\cup W\in\F$ by union closure, and $A'\mapsto A'\cup W$ is injective
on $\F^\emptyset$. (2): for $k\le2$ the coefficient $\phi(S)$ is
nonnegative on $S\in\{\emptyset,W\}$ and the only intermediate level
$|S|=1$ occurs only for $k=2$, where $\phi=0$; combined with (1) the
sum is $\ge0$. (3) isolates the extremes.
\end{proof}

Part~(2) settles Frankl \emph{unconditionally} whenever $\F$ contains a
set of size at most $2$; part~(1) likewise settles the singleton case.
These recover classically known results (\cite{bruhn-schaudt}), but
natively, inside the program's own framework.

\begin{status}
Theorem~\ref{thm:trace-partition} is an exact identity and is correct.
Part~(2) is a genuine unconditional theorem; it reproduces, but does not
improve on, known small-set results. Part~(3) isolates the open case:
the intermediate trace-classes $\{\F^S:0<|S|<k\}$ of a minimal
generator with $k\ge3$.
\end{status}

\subsection{The sub-cube Walsh spectrum and large-fibre localisation (UC3)}
\label{ssec:uc3}

UC3 analyses the intermediate trace-classes via the fibre structure.
Write $U'=U\setminus W$.

\begin{definition}
\label{def:fibre}
For $S\subseteq W$ the \emph{$S$-fibre} is
$\mathcal R_S=\{R\subseteq U':S\sqcup R\in\F\}$, so $|\F^S|=|\mathcal
R_S|$. For $R\subseteq U'$ the \emph{co-fibre} is
$\mathcal S(R)=\{S\subseteq W:S\sqcup R\in\F\}$. The \emph{ambient
fibre} is $\mathcal R_W$.
\end{definition}

\begin{proposition}[UC3 \S2]
\label{prop:fibre-structure}
Each fibre $\mathcal R_S$ and each co-fibre $\mathcal S(R)$ is itself a
union-closed family (on $U'$, respectively on $W$); one has the
graded union laws $\mathcal R_S\vee\mathcal R_{S'}\subseteq\mathcal
R_{S\cup S'}$ and $\mathcal S(R)\vee\mathcal S(R')\subseteq\mathcal
S(R\cup R')$, the universal inclusion $\mathcal R_S\subseteq\mathcal
R_W$ for all $S$, and
\[
  \sum_{x\in W}\bias_x \;=\;\sum_{R\in\mathcal R_W}\Phi\bigl(\mathcal
  S(R)\bigr).
\]
\end{proposition}

The identity reorganises $\sum_{x\in W}\bias_x$ as a sum over the
ambient fibre of the potential of each co-fibre. The co-fibres are
themselves union-closed families on $W$, so --- in the program's phrase
--- the sub-cube Walsh spectrum is ``Frankl about families of
Frankl-families.''

\begin{theorem}[Per-fibre sign theorem and large-fibre localisation, UC3 \S3]
\label{thm:uc3-localisation}
If $|R|\le\lfloor k/2\rfloor$ then $\Phi(\mathcal S(R))\ge0$. Hence the
sign of $\sum_{x\in W}\bias_x$ is governed entirely by the
\emph{large} fibres $\mathcal R_W^{\mathrm{lg}}=\{R:|R|>\lfloor
k/2\rfloor\}$:
\[
  \sum_{x\in W}\bias_x\ge0
  \quad\Longleftarrow\quad
  \sum_{R\in\mathcal R_W^{\mathrm{lg}}}\Phi(\mathcal S(R))
  \;\ge\;-\sum_{R\in\mathcal R_W^{\mathrm{sm}}}\Phi(\mathcal S(R)).
\]
\end{theorem}

\begin{proof}
By minimality of $W$, for $0<|S|<k$ every $R\in\mathcal R_S$ has
$|R|\ge k-|S|$; equivalently, for a small fibre $|R|\le\lfloor
k/2\rfloor$ the co-fibre $\mathcal S(R)$ contains no $S$ with
$0<|S|<\lceil k/2\rceil$, so all its intermediate traces carry
nonnegative potential, and the contributions of $\emptyset$ (if
present) and $W$ are $-k$ and $+k$. The localisation is immediate.
\end{proof}

\begin{corollary}[Bounded generator spread, UC3 \S3]
\label{cor:bounded-spread}
Let $W^+=\bigcup\{A\in\F:A\supseteq W\}$. If $|W^+|\le\lfloor
3k/2\rfloor$ then $\sum_{x\in W}\bias_x\ge0$.
\end{corollary}

This is a genuinely new unconditional region of the open $k\ge3$
frontier: bounded spread forces $\mathcal R_W^{\mathrm{lg}}=\emptyset$.

\begin{remark}
\label{rem:past-sawin}
UC3 argues that the large-fibre object is a \emph{strict refinement} of
the input the entropy method consumes --- it retains the coordinate
coupling that the decoupled entropy relaxation discards --- and is
therefore not subject to the structural reason the entropy method
stalls at $(3-\sqrt5)/2$. We record this carefully: the program does
\emph{not} claim to prove a bound past $(3-\sqrt5)/2$; it claims only to
remove the structural obstruction that prevents the entropy method from
doing so. The distinction is load-bearing and must not be elided.
\end{remark}

\subsection{The trace-defect scalar and the interlock (UC4)}
\label{ssec:uc4}

UC4 isolates the residual into a single scalar. Define the
\emph{cumulative fibre} $\widetilde{\mathcal R}_S=\bigcup_{S_0\subseteq
S}\mathcal R_{S_0}$ with size $q_S=|\widetilde{\mathcal R}_S|$, and the
\emph{trace-defect profile} $d_S=q_S-|\F^S|\ge0$.

\begin{lemma}[Monotone aggregate-sign lemma, UC4 \S1]
\label{lem:monotone-sign}
If $g\colon2^W\to\R$ is monotone non-decreasing then $\langle
g,\phi\rangle=\sum_S g(S)\,\phi(S)\ge0$; if non-increasing then
$\le0$.
\end{lemma}

\begin{proof}
Telescope: $\langle g,\phi\rangle=\sum_{x\in W}\sum_{S\not\ni
x}\bigl(g(S\cup\{x\})-g(S)\bigr)$, each summand $\ge0$.
\end{proof}

\begin{theorem}[Up-closure decomposition and the interlock, UC4 \S2--4]
\label{thm:uc4-interlock}
Write $\sigma(W)=\sum_{x\in W}\bias_x$ and $D(W)=\langle d,\phi\rangle$.
Then $S\mapsto q_S$ is monotone, so by Lemma~\ref{lem:monotone-sign}
the quantity $Q(W):=\langle q,\phi\rangle\ge0$ unconditionally, and
\[
  \sigma(W)\;=\;Q(W)-D(W),\qquad Q(W)\ge0.
\]
Hence $\sigma(W)\ge0$ whenever $D(W)\le Q(W)$, and in particular
whenever $D(W)\le0$. Moreover, ranging $W$ over the set $\mathcal W$ of
all minimal generators,
\[
  \sum_{W\in\mathcal W}\sigma(W)\;=\;\sum_{x}\deg_{\mathcal W}(x)\,
  \bias_x,\qquad \deg_{\mathcal W}(x)=|\{W\in\mathcal W:x\in W\}|,
\]
so $\sum_{W}D(W)\le\sum_{W}Q(W)$ already implies Frankl for $\F$.
\end{theorem}

The identity $\sigma(W)=Q(W)-D(W)$ with $Q(W)\ge0$
\emph{unconditional} compresses the entire residual into the single
trace-defect scalar $D(W)$. The \emph{interlock} (the last display)
relaxes the requirement from a per-generator inequality to a global
one: bad generators may be carried by good ones.

\begin{status}
Theorem~\ref{thm:uc4-interlock} is an exact identity plus a correct
sign on $Q(W)$. It is not a proof of Frankl: $D(W)\le0$ for a single
minimal generator at $k=3$ would already settle Frankl on the open
$m(\F)=3$ frontier, so no single-generator argument can close $D(W)$.
\end{status}

\subsection{Cross-coupling: retraction and Gram form (UC5)}
\label{ssec:uc5}

UC5 quantifies the two couplings the interlock leaves implicit. The
\emph{defect retraction} $\bar\jmath_R(S)=\max(\mathcal S(R)\cap 2^S)$
is the largest trace contained in $S$ still seeing $R$; it is an
interior operator on the up-closure $\mathcal S(R)^\uparrow$, with
fixed-point set $\mathcal S(R)$. This yields an exact closed form for
the cross-trace coupling,
\[
  D(W)\;=\;\sum_{R}\sum_{S_0\in\mathcal S(R)}
  \bigl[(a_{S_0}-1)\,\phi(S_0)+2\,\Sigma(I_{R,S_0})\bigr],
\]
where $I_{R,S_0}$ is an order ideal counting the traces retracting to
$S_0$ and $a_{S_0}=|I_{R,S_0}|$. The cross-generator coupling becomes a
Gram quadratic form: with the generator incidence matrix $M\in
\{0,1\}^{\mathcal W\times U}$ and $G=MM^{\mathsf T}$
($G_{W,W'}=|W\cap W'|$),
\[
  \Sigma_{\mathcal W}\;:=\;\sum_{x}\deg_{\mathcal W}(x)\bigl(2
  \deg_{\mathcal W}(x)-N\bigr)\;=\;2\lVert M^{\mathsf
  T}\mathbf1\rVert^2-kN^2,\qquad N=|\mathcal W|.
\]
UC5 settles two further structured classes --- the
\emph{large-fibre-monotone} class (every large co-fibre is a filter)
and the \emph{concentrated-generator} class ($|U_{\mathcal W}|\le2k$,
via Cauchy--Schwarz on the degree sequence) --- each properly
containing the classes settled in UC3--UC4.

\subsection{The boundary conservation law (UC6)}
\label{ssec:uc6}

UC6 is the terminal note of the combinatorial arc. It converts the
trace-defect into an exact boundary count.

\begin{theorem}[Boundary conservation law, UC6 \S1]
\label{thm:boundary-conservation}
For each large fibre $R$, with the defect $\mathcal
D(R)=\mathcal S(R)^\uparrow\setminus\mathcal S(R)$,
\[
  \Phi(\mathcal D(R))\;=\;\sum_{S_0\in\mathcal S(R)}
  \bigl[a_{S_0}|S_0|-\phi(S_0)\bigr]\;-\;B(R),
\]
where $B(R)\ge|\mathcal S(R)|-1\ge0$ is the number of inter-fibre
Hasse edges of the tiling of $\mathcal S(R)^\uparrow$ by the retraction
fibres. Consequently the large-fibre residual is
\[
  D^{\mathrm{lg}}(W)\;=\;\sum_{R\in\mathcal R_W^{\mathrm{lg}}}
  \Bigl[\sum_{S_0\in\mathcal S(R)}a_{S_0}|S_0|
  -\Phi(\mathcal S(R))-B(R)\Bigr].
\]
\end{theorem}

\begin{proposition}[No per-fibre sign theorem, UC6 \S1]
\label{prop:no-per-fibre}
$\Phi(\mathcal D(R))$ is not sign-definite on the class of large
fibres: both signs occur for union-closure-compatible data. Hence no
theorem of the form ``$\Phi(\mathcal D(R))\le0$ for every large
fibre'' can hold, and the residual is irreducibly a cross-fibre
(sum-level) statement.
\end{proposition}

Proposition~\ref{prop:no-per-fibre} is a genuine negative result: it
proves that the residual \emph{cannot} be closed one fibre at a time.
UC6 also proves a \emph{deficiency-contraction theorem} --- deficiency
is anti-monotone along the fibre lattice $\mathcal R_W$ --- and settles
the $\Gamma$-sparse and low-rooted-shallow classes.

\subsection{Lever-signing on structured classes (UC7--UC8)}
\label{ssec:uc7-uc8}

UC7 and UC8 ``sign the levers'': they take the three exact identities
UC6 produced --- the boundary deficit, the non-generator block, the
excess co-fibre (the failure of the graded union law $\mathcal
S(R)\vee\mathcal S(R')\subseteq\mathcal S(R\cup R')$ to be an
equality) --- and determine the sign of each on progressively larger
structured classes. UC7 works from union closure alone; UC8 admits one
structural input per lever (a cross-fibre persistence invariant; a
cumulative-incidence identity; a square-defect bound). The net output
is a strictly increasing family of structural conditions --- half-level
deficient, $1$-rich, uniformly $W$-balanced and $\Gamma$-sparse,
graded-equal --- each of which implies Frankl, together with worked
witnesses showing the conditions are independent and that none is
implied by union closure alone.

\begin{status}
UC7--UC8 prove genuine implications ``\,structural condition $\Rightarrow$
Frankl for $\F$\,'', and prove (Propositions UC7~4.2, UC8~4.3) that
each condition admits a union-closed family violating it. They enlarge
the settled region; they do not exhaust it. Several separating
witnesses are recorded as structural existence claims rather than fully
written out --- a write-up debt, not an error.
\end{status}

\subsection{Honest terminal status of the combinatorial arc}
\label{ssec:combinatorial-status}

The combinatorial arc is a genuine, monotone chain of reductions. Each
note strictly sharpens its predecessor's residual; no internal
mathematical contradiction was found across UC1--UC8. The arc's honest
terminal state is:

\begin{enumerate}[leftmargin=1.6em]
\item \emph{A complete reduction.} Frankl for $\F$ is equivalent to a
  single explicit integer inequality --- the global sign of the
  boundary deficit $\sum_R B(R)$ against the rooted mass $\sum_R
  \sum_{S_0}a_{S_0}|S_0|$ (Theorem~\ref{thm:boundary-conservation}),
  together with the excess-co-fibre term. Every quantity in it is a
  finite, explicitly computable count.
\item \emph{Genuine unconditional theorems.} The cases $m(\F)\le2$ and
  the singleton case are settled natively
  (Theorem~\ref{thm:trace-partition}); these recover, but do not
  improve, known results.
\item \emph{A growing settled hierarchy.} Bounded generator spread,
  $W$-monotone, low-defect, large-fibre-monotone, concentrated
  generator, $\Gamma$-sparse, low-rooted-shallow, half-level-deficient,
  $1$-rich, uniformly $W$-balanced --- a strictly increasing family of
  structured sub-classes, each settled.
\item \emph{A genuine impossibility result.} No per-fibre sign theorem
  exists (Proposition~\ref{prop:no-per-fibre}); the residual is
  irreducibly global.
\item \emph{Not a proof.} No settled class exhausts union closure; the
  $k=3$ frontier is provably immune to single-generator arguments. The
  final inequality is reduced, explicit, and \emph{open}.
\end{enumerate}

\noindent The combinatorial arc thus delivers a clean, honest
reduction and a real partial-progress hierarchy. It is the conservative
half of the program. The cohomological arc, to which we now turn, is
the half that aimed at a full proof --- and the half where the audit of
Section~\ref{sec:proof-state} finds the live defects.
